Acting on the obligation judgment changes what a bounded view carries. The store's truthful metadata flags $82$ of the $104$ forbidding records; the other $22$ ($18$ retraction-derived, $4$ expiry-derived) are a store-side fact of the kind Table~\ref{tab:believed} records, retrieve as active, and are out of reach of any policy that acts on the metadata. On the $104$ obligation items, the five-entry similarity view contains the forbidding record for $0.5192$; a gold-label fetch of the store's lifecycle records for the obligation's own attribute reaches $0.9519$, $0.4327$ above retrieval alone and short of the $0.50$ bar; gated on the smallest model's committed obligation score at its base rate, the fetch reaches $0.8173$, retaining $0.7561$ of that gain on the supersession-derived items (interval $[0.610, 0.872]$; persona-clustered $[0.562, 0.939]$), short of the $60\%$ bar under the weaker reading. The superseded-names lookup reaches $0.8173$ at the same fetch budget, within tolerance of the model. Naming the retracted and expired attributes raises the gated reach on those obligations only from $0.7805$ to $0.7805$: the gate elicited under the counting summary does not fire on them. The coverage gate at its base rate costs $0.0287$ of coverage on the items whose licensed record the store holds, within tolerance. With the forbidding record in view and tagged, the smallest model gives the stale value as the answer in $0$ of $82$ cases at the registered detector threshold ($11$ of $82$ at the join's gold threshold; $6$ of $17$ where the record printed as active) (Appendix~\ref{app:resp}).
